\documentclass[12pt,letterpaper]{article}

% ─── Paquetes ───────────────────────────────────────────────────────────────
\usepackage[utf8]{inputenc}
\usepackage[T1]{fontenc}
\usepackage[spanish]{babel}
\usepackage[margin=2.5cm]{geometry}
\usepackage{graphicx}
\usepackage{xcolor}
\usepackage{booktabs}
\usepackage{array}
\usepackage{tabularx}
\usepackage{enumitem}
\usepackage{listings}
\usepackage{fancyhdr}
\usepackage{titlesec}
\usepackage{hyperref}
\usepackage{parskip}
\usepackage{amssymb}
\usepackage{pifont}
\usepackage{tcolorbox}
\tcbuselibrary{skins, breakable}

% ─── Colores ────────────────────────────────────────────────────────────────
\definecolor{tecblue}{RGB}{0,70,127}
\definecolor{tecred}{RGB}{200,0,0}
\definecolor{codebg}{RGB}{245,245,245}
\definecolor{codeframe}{RGB}{180,180,180}
\definecolor{keywordcolor}{RGB}{0,0,180}
\definecolor{stringcolor}{RGB}{163,21,21}
\definecolor{commentcolor}{RGB}{0,128,0}
\definecolor{warnbg}{RGB}{255,249,230}
\definecolor{warnframe}{RGB}{220,160,0}
\definecolor{infobg}{RGB}{232,244,255}
\definecolor{infoframe}{RGB}{0,100,200}

% ─── Código fuente ──────────────────────────────────────────────────────────
\lstset{
    backgroundcolor=\color{codebg},
    frame=single,
    rulecolor=\color{codeframe},
    basicstyle=\ttfamily\small,
    keywordstyle=\color{keywordcolor}\bfseries,
    stringstyle=\color{stringcolor},
    commentstyle=\color{commentcolor}\itshape,
    breaklines=true,
    breakatwhitespace=true,
    tabsize=4,
    showstringspaces=false,
    numbers=left,
    numberstyle=\tiny\color{gray},
    numbersep=6pt,
    language=Python,
    literate=
        {á}{{\'a}}1 {é}{{\'e}}1 {í}{{\'i}}1 {ó}{{\'o}}1 {ú}{{\'u}}1
        {Á}{{\'A}}1 {É}{{\'E}}1 {Í}{{\'I}}1 {Ó}{{\'O}}1 {Ú}{{\'U}}1
        {ñ}{{\~n}}1 {Ñ}{{\~N}}1 {¿}{{?`}}1 {¡}{{!`}}1,
}

% ─── Cajas de aviso ─────────────────────────────────────────────────────────
\newtcolorbox{advertencia}{
    colback=warnbg, colframe=warnframe,
    fonttitle=\bfseries, title={\ding{72}~Advertencia},
    breakable
}
\newtcolorbox{info}{
    colback=infobg, colframe=infoframe,
    fonttitle=\bfseries, title={\ding{229}~Nota},
    breakable
}

% ─── Encabezado/pie de página ───────────────────────────────────────────────
\pagestyle{fancy}
\fancyhf{}
\fancyhead[L]{\small CE-5507 Modelación Hardware/Software OO}
\fancyhead[R]{\small Proyecto LinProd — Guía Módulo 3}
\fancyfoot[C]{\thepage}
\renewcommand{\headrulewidth}{0.4pt}

% ─── Secciones ──────────────────────────────────────────────────────────────
\titleformat{\section}{\bfseries\large\color{tecblue}}{}{0em}{}[\titlerule]
\titleformat{\subsection}{\bfseries\normalsize\color{tecblue}}{}{0em}{}
\titleformat{\subsubsection}{\bfseries\normalsize}{}{0em}{}

% ═══════════════════════════════════════════════════════════════════════════
\begin{document}

% ─── Portada ────────────────────────────────────────────────────────────────
\begin{center}
    {\Large\bfseries Proyecto LinProd}\\[0.4em]
    {\LARGE\bfseries\color{tecred} Guía de Implementación — Módulo 3}\\[0.5em]
    {\large\bfseries Interfaz Gráfica de Usuario (GUI)}\\[1em]
    \textit{Instituto Tecnológico de Costa Rica}\\
    \textit{CE-5507 Modelación Hardware/Software Orientado a Objetos}\\
    \textit{I Semestre 2026}
\end{center}

\vspace{0.8em}
\hrule
\vspace{1em}

\begin{info}
Esta guía está dirigida a la persona encargada del \textbf{Módulo 3}. Los
Módulos 1 y 2 ya están completamente implementados y probados; no debes
modificarlos. Tu trabajo consiste exclusivamente en crear
\texttt{linprod/gui/app.py} (y los archivos auxiliares que necesites dentro de
\texttt{linprod/gui/}).
\end{info}

% ─── 1. Contexto general ────────────────────────────────────────────────────
\section{Contexto general del proyecto}

LinProd es un simulador de producción en serie por eventos discretos. El
tiempo avanza en \textbf{ciclos}: en cada ciclo, cada tarea procesa un producto
durante un ciclo de sus \texttt{process\_time} ciclos totales. Los productos
fluyen de tarea en tarea dentro de un proceso, y de proceso en proceso a lo
largo de la línea.

\subsection{Arquitectura de módulos}

\begin{center}
\begin{tabular}{cl}
    \toprule
    Módulo & Responsabilidad \\
    \midrule
    1 & \texttt{linprod/model.py} — clases del dominio (\texttt{Product},
        \texttt{Task}, \texttt{Process}, \texttt{ProductionLine}) \\
    2 & \texttt{linprod/simulator.py} — motor de simulación (\texttt{Simulator}) \\
    \textbf{3} & \textbf{\texttt{linprod/gui/app.py} — GUI (este módulo)} \\
    4 & \texttt{linprod/reports.py} — generación de reportes PDF \\
    \bottomrule
\end{tabular}
\end{center}

La dependencia es unidireccional: el Módulo 3 \emph{importa} de los Módulos 1
y 2, pero éstos no saben nada de la GUI. Nunca modifiques \texttt{model.py} ni
\texttt{simulator.py}.

\subsection{Tecnología}

Usa \textbf{Tkinter} (incluido en Python). Si deseas mejor apariencia visual
puedes usar \textbf{CustomTkinter} (\texttt{pip install customtkinter}), que es
un wrapper moderno sobre Tkinter. La decisión es tuya, pero la GUI debe
ejecutarse con Python 3.8+.

% ─── 2. API de los módulos 1 y 2 ────────────────────────────────────────────
\section{API que debes consumir}

\subsection{Módulo 1 — Construcción de la línea}

Para construir la línea de producción a partir de los datos que el usuario
ingrese en la GUI:

\begin{lstlisting}
from linprod.model import Product, Task, Process, ProductionLine

# Crear un proceso
proceso = Process(name="Ensamblado", is_initial=True)

# Agregar tareas al proceso
proceso.add_task(Task(name="Colocar chasis", process_time=3))
proceso.add_task(Task(name="Instalar motor", process_time=2))

# Construir la linea (add_process enlaza automaticamente next_process)
linea = ProductionLine()
linea.add_process(proceso)
# ... agregar mas procesos ...

# Validar antes de simular (lanza ValueError si hay errores)
linea.validate()
\end{lstlisting}

\subsubsection{Reglas de validación de \texttt{ProductionLine.validate()}}

\begin{itemize}[leftmargin=2em]
    \item Debe haber al menos un proceso.
    \item Exactamente \textbf{un} proceso con \texttt{is\_initial=True}.
    \item Exactamente \textbf{un} proceso con \texttt{is\_final=True}.
    \item Cada proceso debe tener al menos una tarea.
    \item El \texttt{process\_time} de cada tarea debe ser $> 0$.
\end{itemize}

Si \texttt{validate()} lanza \texttt{ValueError}, muestra el mensaje de error
al usuario dentro de la GUI (no dejes que el programa explote sin aviso).

\subsection{Módulo 2 — Control de la simulación}

\begin{lstlisting}
from linprod.simulator import Simulator

# Crear el simulador (linea ya construida y validada)
sim = Simulator(
    line=linea,
    n_products=5,           # cantidad de productos a simular
    on_pause=mi_callback    # funcion que recibe un dict snapshot (opcional)
)

# Iniciar la simulacion (inyecta los productos en el proceso inicial)
sim.start()

# Avanzar UN ciclo (usar esto desde un timer de la GUI)
sim.step()

# Avanzar N ciclos de golpe (util para "run rapido")
sim.run(n_cycles=100)

# Pausar — devuelve snapshot y llama a on_pause si fue registrado
snap = sim.pause()

# Reanudar
sim.resume()

# Reiniciar (mantiene la misma linea, permite cambiar n_products)
sim.reset()

# Propiedades de solo lectura
sim.current_cycle    # int: ciclo actual
sim.is_started       # bool
sim.is_paused        # bool
sim.is_done          # bool: todos los productos completaron la linea

# Estadisticas finales (para el Modulo 4, pero utiles en la GUI tambien)
stats = sim.get_stats()
\end{lstlisting}

\subsection{Snapshot — estructura del diccionario}

\texttt{sim.snapshot()} y el parámetro que recibe \texttt{on\_pause} tienen la
siguiente estructura. Úsala para refrescar la GUI en cada ciclo.

\begin{lstlisting}
{
  "cycle": 7,
  "started": True,
  "paused": False,
  "is_done": False,
  "products_total": 3,
  "products_completed": 1,
  "products_in_progress": 2,

  "line": {
    "process_count": 3,
    "initial_process": "Ensamblado",
    "final_process": "Control de calidad",
    "processes": [
      {
        "process_name": "Ensamblado",
        "is_initial": True,
        "is_final": False,
        "task_count": 2,
        "is_done": False,
        "next_process": "Pintura",
        "tasks": [
          {
            "task_name": "Colocar chasis",
            "process_name": "Ensamblado",
            "process_time": 3,
            "busy": True,
            "current_product_id": 2,   # None si libre
            "remaining_cycles": 1,
            "queue_size": 1,
            "queue_product_ids": [3],
            "products_processed": 1,
            "avg_wait_cycles": 3.0
          },
          # ... mas tareas
        ]
      },
      # ... mas procesos
    ]
  },

  "completed_products": [
    {
      "product_id": 1,
      "entry_time": 0,
      "exit_time": 12,
      "is_completed": True,
      "total_time": 12,
      "task_history": [
        {
          "task_name": "Colocar chasis",
          "process_name": "Ensamblado",
          "start_cycle": 1,
          "end_cycle": 3,
          "wait_cycles": 0
        },
        # ... una entrada por cada tarea que el producto atraveso
      ]
    }
  ],

  "all_products": [ /* mismo formato que completed_products */ ]
}
\end{lstlisting}

% ─── 3. Requerimientos funcionales de la GUI ────────────────────────────────
\section{Requerimientos funcionales que debes cubrir}

Los siguientes requerimientos provienen del enunciado del proyecto. Todos deben
estar presentes en tu implementación.

\subsection{Fase 1 — Configuración de la línea}

\begin{enumerate}[leftmargin=2em]
    \item El usuario puede \textbf{crear procesos} indicando: nombre, si es
          inicial y si es final.
    \item El usuario puede \textbf{agregar tareas} a cada proceso indicando:
          nombre y \texttt{process\_time} (entero $> 0$).
    \item Los procesos deben enlazarse en el orden en que fueron creados para
          formar la línea (\texttt{ProductionLine.add\_process} lo hace
          automáticamente).
    \item El usuario debe poder \textbf{ver la estructura} de la línea que
          está construyendo (lista de procesos con sus tareas).
    \item El usuario debe poder indicar la \textbf{cantidad de productos}
          (\texttt{n\_products}) antes de iniciar.
    \item La GUI debe llamar a \texttt{validate()} antes de iniciar y mostrar
          cualquier error de configuración al usuario de forma clara.
\end{enumerate}

\subsection{Fase 2 — Ejecución y visualización}

\begin{enumerate}[leftmargin=2em]
    \item Botón \textbf{Iniciar}: llama a \texttt{sim.start()} y arranca el
          timer de ciclos.
    \item El \textbf{avance} debe ser visible: la GUI debe actualizarse en
          cada ciclo mostrando el estado de cada tarea (qué producto está
          procesando, cuántos ciclos le restan, tamaño de la cola).
    \item Botón \textbf{Pausar}: llama a \texttt{sim.pause()} y detiene el
          timer. Muestra un ``snapshot'' detallado del estado en ese momento
          (número de ciclo, estado de cada tarea, productos completados,
          productos en proceso, productos en cola).
    \item Botón \textbf{Reanudar}: llama a \texttt{sim.resume()} y retoma el
          timer.
    \item Botón \textbf{Reiniciar}: llama a \texttt{sim.reset()} y regresa a
          la fase de configuración (o al inicio de la simulación con la misma
          línea pero permitiendo cambiar \texttt{n\_products}).
    \item Cuando la simulación termina (\texttt{sim.is\_done == True}), la GUI
          debe notificarlo visualmente y ofrecer la opción de ver el reporte
          (Módulo 4) o reiniciar.
\end{enumerate}

\subsection{Fase 3 — Efectos gráficos}

El enunciado exige \emph{``efectos gráficos''}. Como mínimo debes incluir:

\begin{itemize}[leftmargin=2em]
    \item Indicador visual de que una tarea está \textbf{ocupada} vs.
          \textbf{libre} (p.ej. cambio de color).
    \item Barra de progreso o texto que muestre el avance
          (\texttt{products\_completed / products\_total}).
    \item Resaltar el cuello de botella cuando la simulación termina.
\end{itemize}

Puntos extra (no obligatorio, pero mejora la nota):

\begin{itemize}[leftmargin=2em]
    \item Animación de movimiento del producto entre tareas.
    \item Uso de íconos o colores diferenciados por proceso.
    \item Control de velocidad (ajustar el intervalo del timer).
\end{itemize}

% ─── 4. Implementación sugerida ─────────────────────────────────────────────
\section{Guía de implementación paso a paso}

\subsection{Paso 1 — Estructura de archivos}

Crea los archivos dentro de \texttt{linprod/gui/}. El archivo principal es
\texttt{app.py}; puedes añadir otros si los necesitas.

\begin{lstlisting}[language=bash, numbers=none]
linprod/
  gui/
    __init__.py   (ya existe, no tocar)
    app.py        <- TU TRABAJO
    widgets.py    (opcional: widgets reutilizables)
\end{lstlisting}

\subsection{Paso 2 — Esqueleto de la aplicación}

\begin{lstlisting}
import tkinter as tk
from tkinter import ttk, messagebox
from linprod.model import Task, Process, ProductionLine
from linprod.simulator import Simulator


class LinProdApp(tk.Tk):
    def __init__(self):
        super().__init__()
        self.title("LinProd — Simulador de Producción")
        self.resizable(True, True)

        # Estado de la aplicacion
        self._line: ProductionLine | None = None
        self._sim: Simulator | None = None
        self._timer_id = None            # id del after() activo
        self._cycle_ms = 500             # milisegundos entre ciclos

        # Construir la UI
        self._build_ui()

    # ── Construccion de la interfaz ──────────────────────────────────────────

    def _build_ui(self):
        # Aqui creas los frames, botones, labels, etc.
        pass

    # ── Control de la simulacion ─────────────────────────────────────────────

    def _on_start(self):
        """Llama a sim.start() y arranca el timer."""
        try:
            self._line.validate()
        except ValueError as e:
            messagebox.showerror("Error de configuración", str(e))
            return

        n = int(self._entry_n_products.get())
        self._sim = Simulator(self._line, n_products=n,
                              on_pause=self._on_sim_pause)
        self._sim.start()
        self._schedule_next_cycle()

    def _on_pause(self):
        if self._sim and self._sim.is_started and not self._sim.is_paused:
            if self._timer_id:
                self.after_cancel(self._timer_id)
                self._timer_id = None
            self._sim.pause()   # invoca on_pause automaticamente

    def _on_resume(self):
        if self._sim and self._sim.is_paused:
            self._sim.resume()
            self._schedule_next_cycle()

    def _on_reset(self):
        if self._timer_id:
            self.after_cancel(self._timer_id)
            self._timer_id = None
        if self._sim:
            self._sim.reset()
        # Regresar la GUI a estado inicial
        self._refresh_ui(self._sim.snapshot() if self._sim else None)

    # ── Loop de simulacion (timer Tkinter) ───────────────────────────────────

    def _schedule_next_cycle(self):
        """Programa el siguiente ciclo usando after()."""
        self._timer_id = self.after(self._cycle_ms, self._tick)

    def _tick(self):
        """Avanza un ciclo y refresca la GUI."""
        self._timer_id = None
        if self._sim is None or self._sim.is_paused:
            return

        self._sim.step()
        snap = self._sim.snapshot()
        self._refresh_ui(snap)

        if self._sim.is_done:
            self._on_simulation_done()
        else:
            self._schedule_next_cycle()

    # ── Refresco visual ──────────────────────────────────────────────────────

    def _refresh_ui(self, snap: dict | None):
        """Actualiza todos los widgets con los datos del snapshot."""
        if snap is None:
            return
        # Recorre snap["line"]["processes"] y actualiza labels/colores
        pass

    def _on_sim_pause(self, snap: dict):
        """Callback que llama el Simulator cuando se pausa."""
        self._refresh_ui(snap)
        # Mostrar ventana/panel de estado detallado
        pass

    def _on_simulation_done(self):
        """Acciones al completar todos los productos."""
        messagebox.showinfo(
            "Simulación finalizada",
            f"Todos los productos completaron la línea en "
            f"{self._sim.current_cycle} ciclos."
        )


def main():
    app = LinProdApp()
    app.mainloop()


if __name__ == "__main__":
    main()
\end{lstlisting}

\begin{info}
El método clave para el loop de animación es \texttt{after()}. Nunca uses
\texttt{time.sleep()} dentro de Tkinter porque bloquea la interfaz. Con
\texttt{after(ms, funcion)} programas que Tkinter llame a \texttt{funcion}
después de \texttt{ms} milisegundos sin bloquear nada.
\end{info}

\subsection{Paso 3 — Sección de configuración de la línea}

Esta sección debe permitir al usuario construir la línea \textbf{antes} de
iniciar la simulación. Una estructura sugerida:

\begin{enumerate}[leftmargin=2em]
    \item Un \texttt{Frame} izquierdo con el formulario para agregar
          proceso/tarea.
    \item Un \texttt{Listbox} o \texttt{Treeview} que muestre la estructura
          actual de la línea.
    \item Botones: \textbf{Agregar Proceso}, \textbf{Agregar Tarea},
          \textbf{Eliminar} (opcional), \textbf{Limpiar línea}.
\end{enumerate}

\begin{lstlisting}
def _build_config_frame(self, parent):
    frame = ttk.LabelFrame(parent, text="Configuración de la línea")

    # Nombre del proceso
    ttk.Label(frame, text="Nombre del proceso:").grid(row=0, column=0)
    self._entry_proc_name = ttk.Entry(frame)
    self._entry_proc_name.grid(row=0, column=1)

    # Checkboxes inicial / final
    self._var_is_initial = tk.BooleanVar()
    self._var_is_final   = tk.BooleanVar()
    ttk.Checkbutton(frame, text="Inicial",
                    variable=self._var_is_initial).grid(row=1, column=0)
    ttk.Checkbutton(frame, text="Final",
                    variable=self._var_is_final).grid(row=1, column=1)

    ttk.Button(frame, text="+ Agregar proceso",
               command=self._add_process).grid(row=2, column=0, columnspan=2)

    # Separador
    ttk.Separator(frame, orient="horizontal").grid(
        row=3, column=0, columnspan=2, sticky="ew", pady=6)

    # Nombre de la tarea y tiempo
    ttk.Label(frame, text="Nombre de la tarea:").grid(row=4, column=0)
    self._entry_task_name = ttk.Entry(frame)
    self._entry_task_name.grid(row=4, column=1)

    ttk.Label(frame, text="Tiempo de procesamiento:").grid(row=5, column=0)
    self._entry_task_time = ttk.Entry(frame)
    self._entry_task_time.grid(row=5, column=1)

    ttk.Button(frame, text="+ Agregar tarea",
               command=self._add_task).grid(row=6, column=0, columnspan=2)

    return frame

def _add_process(self):
    name = self._entry_proc_name.get().strip()
    if not name:
        messagebox.showwarning("Dato faltante", "Ingrese el nombre del proceso.")
        return
    proc = Process(name,
                   is_initial=self._var_is_initial.get(),
                   is_final=self._var_is_final.get())
    if self._line is None:
        self._line = ProductionLine()
    self._line.add_process(proc)
    self._refresh_line_tree()

def _add_task(self):
    # Requiere que haya al menos un proceso creado
    if self._line is None or not self._line.processes:
        messagebox.showwarning("Sin proceso", "Cree primero un proceso.")
        return
    name = self._entry_task_name.get().strip()
    try:
        t = int(self._entry_task_time.get())
        if t <= 0:
            raise ValueError
    except ValueError:
        messagebox.showerror("Valor inválido",
                             "El tiempo debe ser un entero > 0.")
        return
    # La tarea se agrega al ULTIMO proceso creado
    last_process = self._line.processes[-1]
    last_process.add_task(Task(name, process_time=t))
    self._refresh_line_tree()
\end{lstlisting}

\subsection{Paso 4 — Panel de visualización de la simulación}

Para mostrar el estado de cada tarea en tiempo real, un \texttt{Treeview} o
una cuadrícula de \texttt{Frame}/\texttt{Label} funciona bien. La clave es
iterar sobre \texttt{snap["line"]["processes"]} en \texttt{\_refresh\_ui}.

\begin{lstlisting}
def _refresh_ui(self, snap: dict | None):
    if snap is None:
        return

    self._lbl_cycle.config(text=f"Ciclo: {snap['cycle']}")
    self._lbl_completed.config(
        text=f"Completados: {snap['products_completed']} / "
             f"{snap['products_total']}"
    )

    for proc_data in snap["line"]["processes"]:
        for task_data in proc_data["tasks"]:
            key = task_data["task_name"]
            if key in self._task_widgets:
                widget = self._task_widgets[key]
                if task_data["busy"]:
                    pid = task_data["current_product_id"]
                    rem = task_data["remaining_cycles"]
                    widget.config(
                        text=f"P{pid} ({rem} ciclos restantes)",
                        background="#ffcc66"   # amarillo = ocupada
                    )
                else:
                    widget.config(
                        text="Libre",
                        background="#99dd99"   # verde = libre
                    )
\end{lstlisting}

\subsection{Paso 5 — Control de velocidad (opcional pero recomendado)}

Un \texttt{Scale} (slider) que modifique \texttt{\_cycle\_ms} permite al
usuario acelerar o ralentizar la simulación:

\begin{lstlisting}
def _build_speed_control(self, parent):
    ttk.Label(parent, text="Velocidad (ms/ciclo):").pack(side="left")
    self._scale_speed = ttk.Scale(
        parent, from_=50, to=2000,
        orient="horizontal",
        command=lambda v: setattr(self, "_cycle_ms", int(float(v)))
    )
    self._scale_speed.set(500)
    self._scale_speed.pack(side="left", fill="x", expand=True)
\end{lstlisting}

\subsection{Paso 6 — Ventana de snapshot al pausar}

Cuando el usuario pausa, el Simulator llama a \texttt{on\_pause(snap)}. Usa
esa oportunidad para mostrar una ventana \texttt{Toplevel} con el detalle:

\begin{lstlisting}
def _on_sim_pause(self, snap: dict):
    self._refresh_ui(snap)
    win = tk.Toplevel(self)
    win.title(f"Estado — Ciclo {snap['cycle']}")

    text = tk.Text(win, width=60, height=30, font=("Courier", 10))
    text.pack(fill="both", expand=True)

    text.insert("end", f"=== SNAPSHOT — Ciclo {snap['cycle']} ===\n\n")
    text.insert("end",
        f"Productos completados : {snap['products_completed']}\n"
        f"Productos en proceso  : {snap['products_in_progress']}\n\n"
    )
    for proc in snap["line"]["processes"]:
        text.insert("end", f"Proceso: {proc['process_name']}\n")
        for task in proc["tasks"]:
            estado = (
                f"procesando P{task['current_product_id']} "
                f"({task['remaining_cycles']} ciclos restantes)"
                if task["busy"] else "LIBRE"
            )
            text.insert("end",
                f"  └ {task['task_name']:25s} | {estado} | "
                f"cola: {task['queue_size']}\n"
            )
        text.insert("end", "\n")

    text.config(state="disabled")
\end{lstlisting}

% ─── 5. Integración con el Módulo 4 ─────────────────────────────────────────
\section{Integración con el Módulo 4 (Reportería)}

Cuando la simulación termina, la GUI debe ofrecer un botón \textbf{Generar
reporte}. La interfaz entre el Módulo 3 y el Módulo 4 es la siguiente:

\begin{lstlisting}
from linprod.reports import generate_report   # Modulo 4

def _on_generate_report(self):
    if self._sim and self._sim.is_done:
        stats = self._sim.get_stats()
        generate_report(stats)   # el Modulo 4 define esta funcion
\end{lstlisting}

Mientras el Módulo 4 no esté listo, puedes dejar el botón desactivado o hacer
que muestre los stats en una ventana de texto como las del paso 6.

% ─── 6. Checklist de entrega ────────────────────────────────────────────────
\section{Checklist antes de entregar}

\begin{itemize}[leftmargin=2em, label=\( \square \)]
    \item La GUI arranca sin errores con \texttt{python3 -m linprod.gui.app}.
    \item Se pueden agregar procesos y tareas en tiempo de ejecución (sin
          límite).
    \item El proceso inicial y final se marcan explícitamente.
    \item \texttt{validate()} es llamado antes de iniciar; sus errores se
          muestran al usuario.
    \item Los botones Iniciar / Pausar / Reanudar / Reiniciar funcionan.
    \item Al pausar se muestra el estado detallado de cada tarea.
    \item El estado de cada tarea se actualiza visualmente en cada ciclo.
    \item Hay al menos un efecto gráfico (colores según estado busy/libre,
          barra de progreso, etc.).
    \item Al terminar se notifica al usuario con el ciclo final.
    \item El código no modifica \texttt{model.py} ni \texttt{simulator.py}.
    \item Los imports son desde \texttt{linprod.model} y
          \texttt{linprod.simulator}.
\end{itemize}

% ─── 7. Advertencias frecuentes ─────────────────────────────────────────────
\section{Errores comunes a evitar}

\begin{advertencia}
\textbf{Nunca uses \texttt{time.sleep()} dentro de Tkinter.} Bloquea el hilo
principal y la ventana deja de responder. Usa siempre \texttt{after()}.
\end{advertencia}

\begin{advertencia}
\textbf{Nunca llames a \texttt{sim.step()} en un hilo aparte.} Tkinter y los
objetos de los Módulos 1/2 no son thread-safe. Todo el loop de simulación debe
ocurrir en el hilo principal mediante \texttt{after()}.
\end{advertencia}

\begin{advertencia}
\textbf{No importes desde \texttt{linprod.gui} dentro de \texttt{model.py} ni
\texttt{simulator.py}.} La dependencia es unidireccional: GUI $\rightarrow$
Simulator $\rightarrow$ Model.
\end{advertencia}

\begin{info}
Si usas \textbf{CustomTkinter}, todos los principios son los mismos. Solo
reemplaza \texttt{tk.Tk} por \texttt{ctk.CTk}, \texttt{ttk.Button} por
\texttt{ctk.CTkButton}, etc. El \texttt{after()} y el \texttt{mainloop()} son
idénticos.
\end{info}

% ─── 8. Cómo ejecutar ────────────────────────────────────────────────────────
\section{Cómo ejecutar la aplicación}

Desde la raíz del repositorio (\texttt{LinProd-ModelacionHardware/}):

\begin{lstlisting}[language=bash, numbers=none]
# Ejecutar la GUI
python3 -m linprod.gui.app

# Ejecutar los tests unitarios (cuando esten escritos)
python3 -m pytest tests/

# Verificar que el Modulo 1 sigue funcionando
python3 linprod/model.py

# Verificar que el Modulo 2 sigue funcionando
python3 -m linprod.simulator
\end{lstlisting}

\begin{info}
El Módulo 2 (\texttt{simulator.py}) debe ejecutarse como módulo
(\texttt{python3 -m linprod.simulator}), no como script directo, porque
utiliza imports absolutos del paquete \texttt{linprod}.
\end{info}

\end{document}
